\section{Introduction}
\label{Introduction}


Extracting building footprints has long been an essential task in aerial imagery interpretation and GIScience, supporting applications such as disaster management, urban planning, and infrastructure monitoring \citep{Awrangjeb_AutomaticExtraction}. Accurate and up-to-date building footprints are not only required by national mapping agencies for maintaining topographic databases, but also play a critical role in crisis situations \citep{ZHENG_Buildingdamage}. However, vectorizing building footprints over large areas requires substantial manual effort. Mapping agencies must invest significant man-hours into digitizing and validating footprints, which naturally motivates the development of automated or semi-automated approaches \citep{rastogi2022automatic}.

Over the past decades, many different methods have been proposed \citep{ZHENG_Buildingdamage}, yet no approach has achieved 100 percent accuracy \citep{rastogi2022automatic}, largely due to the vast diversity of building shapes, materials, and architectural concepts.

These challenges have increasingly motivated the use of artificial intelligence and deep learning methods for automated building-footprint extraction. The topics of Artificial Intelligence, Machine Learning, Deep Learning, and Neural Networks are closely related and build on one another \citep{shiri2023comprehensive}. They can be understood as a nested hierarchy: neural networks form the foundation; machine learning builds upon them; deep learning represents a layered subset of machine learning; and artificial intelligence encompasses all of these concepts. In this work, the focus lies specifically on deep learning using pretrained transformer-based
foundation models.

With the recent development of artificial intelligence, researchers have increasingly explored foundation models instead of purely task-specific deep-learning architectures. The Vision Transformer (ViT) \citep{DosovitskiyTransformer} was among the first transformer-based models for vision tasks, demonstrating that attention mechanisms can effectively be applied to image recognition. Building on this idea, the Dense Prediction Transformer (DPT) \citep{ranftl2021visiontransformersdenseprediction} showed that transformers can also produce coherent global predictions for dense visual tasks.

More recently, research has moved toward integrating visual foundation models with multimodal large language models (MLLMs) \citep{wu2023multimodallargelanguagemodels}. These models combine visual and textual reasoning, offering new possibilities for interpreting spatial data. In GIScience, such approaches could enhance spatial reasoning and enable automated analysis of geospatial features beyond pixel-level classification \citep{Ji_2025_foundationGeospatiolreasoning}. Early studies have begun to explore such combinations, typically allowing the language model to guide, describe, or evaluate the visual model's segmentation process. This development is directly relevant for this thesis, where MLLM-based visual assessment is combined with SAM-based segmentation for selective building-footprint refinement.

This thesis takes a different approach. Instead of using a language model to directly control segmentation, a two-step framework is proposed in which the MLLM evaluates existing building polygons and supports selective refinement, thereby clearly separating assessment and correction. This follows an update-oriented geospatial workflow, where the aim is not to reconstruct an entire building dataset from scratch, but to identify and improve features that require correction \citep{Qin_changedetection}.

In the first stage, imperfect building polygons from Google Open Buildings serve as the starting point. These polygons are evaluated by a multimodal large language model, which identifies whether a building is visible, directs the case into the appropriate refinement path, and generates spatial prompts where required. In the second stage, a visual foundation model generates refined geometry candidates, which are subsequently converted and regularized into GIS-ready building polygons.

This design addresses key limitations of current end-to-end approaches, which typically require large labeled datasets and are difficult to adapt across regions or data domains. By decoupling assessment, segmentation, and correction, the proposed workflow enables clearer feedback loops and greater flexibility in integrating automated quality assurance into geospatial production processes.

The research gap addressed in this thesis lies in the limited transfer of such multimodal, feedback-driven AI approaches to GIScience, particularly for the automated quality assessment of vector-based building footprints.


\textbf{RQ1:} To what extent can a multimodal large language model support the visual quality assessment of existing vector-based building footprints?

\textbf{RQ2:} How effectively can MLLM-generated spatial prompts, combined with SAM-based segmentation and geometric post-processing, improve selected building-footprint geometries?

\textbf{RQ3:} How transferable is the proposed workflow across different geographic regions, building morphologies, and imagery conditions?
